% vim: tw=80 noai
\documentclass[normaltoc,capchap,capsec,times]{abnt}
\usepackage[latin1]{inputenc}
\usepackage[T1]{fontenc}
\usepackage[brazil]{babel}
\usepackage[alf]{abntcite}
\usepackage[ordem=alf]{tabela-simbolos}
\usepackage{url}
\usepackage{graphicx}
\usepackage{listings}
\usepackage{verbatim}
\usepackage{subfigure}
\usepackage{multicol}
\usepackage{framed}
% O comando abaixo define o diretorio onde devem ser colocadas as imagens.
% Neste caso o diretorio e' ./imagens/
\graphicspath{ {./imagens/} }
\def\lstlistingname{Listagem}

%%%%%%%%%%%%%%%%%%%%%%%%%%%%%%%%%%%%%%%%%%%%%%%%%%%%
% Dados da monografia
%%%%%%%%%%%%%%%%%%%%%%%%%%%%%%%%%%%%%%%%%%%%%%%%%%%%

\newcommand{\meunome}{Fulano de tal}
\newcommand{\meutitulo}{Um modelo LaTeX para monografias no DCC-UFBA}
\newcommand{\meusubtitulo}{Uma abordagem pr�tica}
\newcommand{\meuano}{2013}
\newcommand{\meuorientador}{Orientadora: \profa\ Christina von Flach Garcia
Chavez}

%%%%%%%%%%%%%%%%%%%%%%%%%%%%%%%%%%%%%%%%%%%%%%%%%%%%

%% O comando \obs aqui definido permite que o autor faca anotacoes na
%% monografia que aparecem no PDF gerado. Para ativar o comando, descomente
%% a primeira linha e comente a segunda.
%% Exemplo de uso: \obs{Preciso melhorar este par�grafo...}

%\newcommand{\obs}[1]{\underline{\textbf{OBSERVA��O}}: \emph{#1}}
\newcommand{\obs}[1]{}

\def\ordfem{\mbox{\raise .35em \hbox{\underline{\scriptsize a}\ }}}
\def\ordmasc{\mbox{\raise .35em \hbox{\underline{\scriptsize o}\ }}}
\def\profa{Prof\ordfem.}

%%%%%%%%%%%%%%%%%%%%%%%%%%%%%%%%%%%%%%%%%%%%%%%%%%%%

\begin{document}

% Capa com Bras�o

\begin{titlepage}
   \begin{center}
    	%logotipo
               \includegraphics[scale=0.3]{brasao_ufba} \\
	%\vspace{0.7in}
              \centering{ 
	      \bf{
	      \LARGE{
		\uppercase{UNIVERSIDADE FEDERAL DA BAHIA} \\
 	      }
	      \Large {
                   	\uppercase{INSTITUTO DE MATEM�TICA} \\
	      }
                   \large {
                       \uppercase{DEPARTAMENTO DE CI�NCIA DA COMPUTA��O} \\
                  }
              } }
   \end{center}
   \vfill
   \begin{center}
       \bf{
       \large{\uppercase{\meunome}  \\  }
       }
   \end{center}
   \vspace{0.2in}
   \begin{center}
       \bf{
      	 \LARGE{ \uppercase{\meutitulo} } \\
      	 \Large{ \uppercase{\meusubtitulo} }
         \obs{\\ \Large{Esta vers�o da monografia cont�m coment�rios do autor.
          Para remov�-los, redefina o comando LaTeX \texttt{obs}.}}
       }
   \end{center}

   \vfill
   \hspace{\stretch{1}}
   \vfill
   \begin{center}
      \normalsize{
          Salvador \\
          \meuano
       }
   \end{center}

\end{titlepage}

%comando abaixo cria uma capa redundante, mas como a capa com bras�o foi 
% feita 'manualmente', n�o faz sentido usar este comando:
%\capa



%\folhaderosto
% o comando acima foi comentado para n�o criar uma folha de 
% rosto redundante, j� que ela feita 'manualmente' abaixo

\begin{titlepage}
 \vfill
 \begin{center}
   {\large \uppercase{ \bf{ \meunome\ } } } \\[7cm]
   {\Huge \uppercase{ \bf{ \meutitulo\ } } }\\[1cm]
   \vfill
   \hspace{.45\textwidth} % posicionando a minipage
   \begin{minipage}{.5\textwidth}
     \begin{espacosimples}
       \bf{
	Monografia apresentada ao Curso de gradua��o em Ci�ncia da Computa��o, 
	Departamento de Ci�ncia da Computa��o, Instituto de Matem�tica,  Universidade Federal da 
	Bahia, como requisito parcial para obten��o do grau de Bacharel em  Ci�ncia da Computa��o. \\ 
       }      
     \end{espacosimples}
     \begin{espacosimples}    
       \meuorientador
     \end{espacosimples}
   \end{minipage}
   \vfill
   Salvador \\
   \meuano
 \end{center}
\end{titlepage}


\begin{resumo}
Nonono nonono nonono, nonono, nonono nonono nonono nononono nonno. 
Nonono nonono nonono, nonono, nonono nonono nonono nononono nonno. 
Nonono nonono nonono, nonono, nonono nonono nonono nononono nonno. 
Nonono nonono nonono, nonono, nonono nonono nonono nononono nonno. 
Nonono nonono nonono, nonono, nonono nonono nonono nononono nonno. 
Nonono nonono nonono, nonono, nonono nonono nonono nononono nonno. 
Nonono nonono nonono, nonono, nonono nonono nonono nononono nonno. 

\textbf{Palavras-chave:}
monografia,
gradua��o, 
projeto final.
\end{resumo}

% O abstract e' opcional.
\begin{abstract}
Nonono nonono nonono, nonono, nonono nonono nonono nononono nonno. 
Nonono nonono nonono, nonono, nonono nonono nonono nononono nonno. 
Nonono nonono nonono, nonono, nonono nonono nonono nononono nonno. 
Nonono nonono nonono, nonono, nonono nonono nonono nononono nonno. 
Nonono nonono nonono, nonono, nonono nonono nonono nononono nonno. 
Nonono nonono nonono, nonono, nonono nonono nonono nononono nonno. 
Nonono nonono nonono, nonono, nonono nonono nonono nononono nonno. 

\textbf{Keywords:} 
monograph,
graduation,
final project.
\end{abstract}

%% As listas a seguir sao opcionais:
\listadefiguras
%\listadetabelas
\listadesiglas
%\listadesimbolos

\sumario

% O conteudo da monografia esta' nos seguintes arquivos:
\chapter{Introdu��o}

Modelo de monografia usando as normas ABNT (Associa��o Brasileira de Normas
T�cnicas) \sigla{ABNT}{Associa��o Brasileira de Normas T�cnicas} e adapta��o
personalizada do padr�o do Departamento de Ci�ncia da Computa��o (DCC) da
Universidade Federal da Bahia (UFBA).
\sigla{DCC}{Departamento de Ci�ncia da Computa��o}
\sigla{UFBA}{Universidade Federal da Bahia}
Fontes latex cedidos pela ABNT e disponibilizados por 
Maur�cio Vieira. (Valeu Maurix!). Adaptado por Abelmon Bastos por solicita��o
da \profa\ D�bora Abdalla para o semestre 2005.1.
Adaptado por Rodrigo Rocha por solicita��o da \profa\ D�bora Abdalla no fim
do semestre 2007.1.
Atualizado por Ibirisol em 2013.1 para adequa��o ao \cite{manualVisual} da UFBA.

% use sua propria estrutura

%Problema etc

%Objetivo etc

%Resultados esperados etc

%Estrutura/Organiza��o etc

\chapter{Rede sem fio em malha} \label{cap:WMN}

Este trabalho relata o onon ono non ono non ono non ono non ono non on.

\section{Listagens} \label{sec:listagens}

Na listagem \ref{lst:testjUnit} 
� mostrado o teste do m�todo \texttt{Engine.initialize()}:

\lstset{language=java}
\lstset{commentstyle=\textit}
\begin{lstlisting}[frame=trbl, caption=Classe Factory2D,label=lst:testjUnit]{}
public class EngineTest
// JUnitDoclet begin extends_implements
extends TestCase
// JUnitDoclet end extends_implements
{
  //...
  public void testInitialize() throws Exception {
   // JUnitDoclet begin method initialize
   EngineState engineState = (EngineState) PrivateAccessor.
    getField(engine,"engineState");
   engine.initialize();
   assertEquals(engineState, new InitState());
   // JUnitDoclet end method initialize
  }
  ...
}
\end{lstlisting}

\section{Figuras} \label{sec:figuras}

� poss�vel usar imagens vetoriais no formato PDF, como pode ser visto
na figura \ref{fig:ufba}, ou imagens \emph{bitmap} no formato PNG, como
a da figura \ref{fig:ufba2}.

\begin{figure}
\centering
\includegraphics[scale=0.3]{brasao_ufba_vetorial}
\caption{Bras�o da UFBA (vetorial)}
\label{fig:ufba}
\end{figure}

\begin{figure}
\centering
\includegraphics[width=0.3\textwidth]{brasao_ufba}
\caption{Bras�o da UFBA (\emph{bitmap})}
\label{fig:ufba2}
\end{figure}


\chapter{Outro cap�tulo}
\label{cap:outrocapitulo}

Como visto no cap�tulo \ref{cap:conceitos}, no nonno\footnote{Isto � uma nota
de rodap�.} no nonon onono:
\begin{itemize}
  \item{nononoo}
  \item{nononono}
  \item{no}
\end{itemize}


Nonono nonnono onononono \cite{fowler2000}.

\chapter{Conceitos} \label{cap:conceitos}

Este trabalho relata o onon ono non ono non ono non ono non ono non on.

\section{Listagens} \label{sec:listagens}

Na listagem \ref{lst:testjUnit} 
� mostrado o teste do m�todo \texttt{Engine.initialize()}:

\lstset{language=java}
\lstset{commentstyle=\textit}
\begin{lstlisting}[frame=trbl, caption=Classe Factory2D,label=lst:testjUnit]{}
public class EngineTest
// JUnitDoclet begin extends_implements
extends TestCase
// JUnitDoclet end extends_implements
{
  //...
  public void testInitialize() throws Exception {
   // JUnitDoclet begin method initialize
   EngineState engineState = (EngineState) PrivateAccessor.
    getField(engine,"engineState");
   engine.initialize();
   assertEquals(engineState, new InitState());
   // JUnitDoclet end method initialize
  }
  ...
}
\end{lstlisting}

\section{Figuras} \label{sec:figuras}

� poss�vel usar imagens vetoriais no formato PDF, como pode ser visto
na figura \ref{fig:ufba}, ou imagens \emph{bitmap} no formato PNG, como
a da figura \ref{fig:ufba2}.

\begin{figure}
\centering
\includegraphics[scale=0.3]{brasao_ufba_vetorial}
\caption{Bras�o da UFBA (vetorial)}
\label{fig:ufba}
\end{figure}

\begin{figure}
\centering
\includegraphics[width=0.3\textwidth]{brasao_ufba}
\caption{Bras�o da UFBA (\emph{bitmap})}
\label{fig:ufba2}
\end{figure}


\chapter{Outro cap�tulo}
\label{cap:outrocapitulo}

Como visto no cap�tulo \ref{cap:conceitos}, no nonno\footnote{Isto � uma nota
de rodap�.} no nonon onono:
\begin{itemize}
  \item{nononoo}
  \item{nononono}
  \item{no}
\end{itemize}


Nonono nonnono onononono \cite{fowler2000}.

\chapter{Conclus�o}

O trabalho onon ono non ono non ono non ono non ono non on
n ono non ono non ono ono non ononon ono non ono non ono non ono non on
 
%\section{Pontos fortes da metodologia proposta}

\section{Dificuldades encontradas}

O trabalho onon ono non ono non ono non ono non ono non on
n ono non ono non ono ono non ononon ono non ono non ono non ono non on

\section{Trabalhos futuros}

Pode-se indicar como trabalhos futuros:

\textbf{n ono non ono non ono non ono non }.
n ono non ono non ono non ono non n ono non 

ono non ono non ono non n ono non ono non ono non ono non 
\textbf{controlador} n ono non ono non ono ono non ononon ono non ono non ono non ono non on
n ono no oo non ono ononon ono  non ono non ono non ono non ono non on

\textbf{ono non ono}
o non ono non ono ono non ononon ono non ono non ono non ono non on
n ono no oo non ono ononon ono  non ono non ono non ono non ononon o 

%\section{Divulga��o}

%\section{Pontos fortes e fracos}


\appendix

\chapter{Resultados experimentais}

No no nnononono no n ono o nn.


\bibliography{monografia}

\end{document}
